\section{Introduction}
Distribution feeder reconfiguration changes which sectionalizing and tie
switches conduct, allowing the same physical infrastructure to supply loads
through different connected configurations. Classical work treats this as an
electrical-network problem with objectives such as loss reduction and load
balancing \cite{baran1989}. Software-defined networking (SDN), by comparison,
separates programmable forwarding decisions from the mechanisms that execute
them in a communication network \cite{mckeown2008}. The conceptual separation
is attractive for a project in which one team develops network-control logic
and another investigates physical low-voltage (LV) switching.

The analogy has an important boundary. Alternating-current power is not
forwarded as independently addressed packets through passive feeder branches.
Switching alters circuit connectivity, after which voltages and currents follow
the network impedances, source characteristics, and electrical dynamics.
Consequently, a route returned by a networking algorithm must be translated
into a physically represented switching sequence and evaluated in an
electrical simulator. A graph can establish connectivity; it cannot, by
itself, establish acceptable voltage, equipment loading, transient behavior,
or protection coordination.

The broader project aims to maintain a communication control plane independently
of switchable power paths. Its present electrical realization is a single-source
LV feeder, not a complete multi-resource microgrid. This distinction matters:
the U.S. Department of Energy's microgrid description includes local energy
resources, controllability as an entity, and grid-connected or islanded
operation \cite{doeMicrogrids}. The current case has neither distributed
energy resource (DER) control nor a demonstrated islanding transition. The
term ``toward'' in the title therefore denotes a research trajectory rather
than an already validated operating capability.

This paper asks three practical questions. First, can externally specified
network commands be mapped consistently onto physical breaker schedules while
retaining identifiers and operation-level decisions? Second, do the saved EMT
measurements show the expected distinction between fault inception, isolation,
and restoration of healthy loads? Third, can graphical dashboard edits trigger
fresh, inspectable PSCAD experiments without changing the networking team's
implementation?

The resulting contributions are:
\begin{enumerate}
\item A modular command-to-power-plane workflow that retains a structured
controller contract, operation-level graph checks, explicit component bindings,
and a reproducible timed experiment manifest.
\item A PSCAD automation path with native breaker and fault components,
serialized run ownership, fresh-result checks, and a defined set of 31 research
channels exposed through a graphical JSON-backed dashboard.
\item An evidence-bounded analysis of six saved deterministic scenarios,
including a physical fault, downstream restoration, and an explicitly imposed
restoration delay, accompanied by reproducible publication assets.
\end{enumerate}
These are integration and functional-validation contributions. No superiority
over established reconfiguration algorithms or cyber-physical testbeds is
claimed. The experiments support further research on the interface between
network decisions and electrical consequences; they are not evidence that
the proposed architecture is ready to operate a physical feeder.
